\begin{figure}[htbp]
	\centering
	%\resizebox{\textwidth}{!}{%
		\begin{tikzpicture}[>=Stealth, line width=0.8pt, scale=0.85]
			
			% ═══════════════════════════════════════════════════════════
			% PANNELLO (a) — Prestazione cinematica del bipendolo
			% (realizzazione: pantografo a parallelogramma)
			% ═══════════════════════════════════════════════════════════
			\begin{scope}[xshift=0cm]
				
				%% ── Parametri ──
				\def\xAuno{-0.40}
				\def\xAdue{ 0.40}
				\def\yA{ 0.00}
				\def\Lbiella{0.70}      % era 0.90
				\def\angSup{-110}
				\def\angInf{-70}
				\def\rn{0.10}
				\def\tw{0.22}
				\def\thA{0.32}
				\def\gw{0.40}
				
				% Coordinate calcolate
				\pgfmathsetmacro{\dxSup}{\Lbiella*cos(\angSup)}
				\pgfmathsetmacro{\dySup}{\Lbiella*sin(\angSup)}
				\pgfmathsetmacro{\dxInf}{\Lbiella*cos(\angInf)}
				\pgfmathsetmacro{\dyInf}{\Lbiella*sin(\angInf)}
				% B_1, B_2 (nodi traversa)
				\pgfmathsetmacro{\xBuno}{\xAuno+\dxSup}
				\pgfmathsetmacro{\yB}{\yA+\dySup}
				\pgfmathsetmacro{\xBdue}{\xAdue+\dxSup}
				% G_1, G_2 (cerniere a terra)
				\pgfmathsetmacro{\xGuno}{\xBuno+\dxInf}
				\pgfmathsetmacro{\yG}{\yB+\dyInf}
				\pgfmathsetmacro{\xGdue}{\xBdue+\dxInf}
				
				
				%% ── Corpo C (più grande, agganciato a A_1=-0.40 e A_2=+0.40) ──
				\draw[line width=1pt] 
				plot[smooth cycle, tension=0.7] coordinates 
				{(-0.40, 0.00) (-1.30, 0.35) (-1.80, 1.15) 
					(-1.45, 2.05) (-0.55, 2.40) (0.55, 2.40) 
					(1.45, 2.05) (1.80, 1.15) (1.30, 0.35) (0.40, 0.00)};
				\node at (-0.65, 1.65) {$\mathcal{C}$};
				\node[font=\footnotesize] at (0.00, 0.75) {$\vartheta = 0$};
				
				%% ── Bielle superiori (parallelogramma 1) ──
				\draw[line width=0.8pt] (\xAuno, \yA) -- (\xBuno, \yB);
				\draw[line width=0.8pt] (\xAdue, \yA) -- (\xBdue, \yB);
				
				%% ── Traversa B_1 B_2 ──
				\draw[line width=0.8pt] (\xBuno, \yB) -- (\xBdue, \yB);
				
				%% ── Bielle inferiori (parallelogramma 2) ──
				\draw[line width=0.8pt] (\xBuno, \yB) -- (\xGuno, \yG);
				\draw[line width=0.8pt] (\xBdue, \yB) -- (\xGdue, \yG);
				
				%% ── Cerniere al suolo G_1, G_2 ──
				\draw (\xGuno, \yG) 
				-- ++(-\tw, -\thA) -- ++(2*\tw, 0) -- cycle;
				\draw (\xGuno-\gw, \yG-\thA) -- (\xGuno+\gw, \yG-\thA);
				\foreach \x in {-0.30,-0.15,0,0.15,0.30}
				\draw[thin] (\xGuno+\x, \yG-\thA) -- ++(-0.10, -0.13);
				
				\draw (\xGdue, \yG) 
				-- ++(-\tw, -\thA) -- ++(2*\tw, 0) -- cycle;
				\draw (\xGdue-\gw, \yG-\thA) -- (\xGdue+\gw, \yG-\thA);
				\foreach \x in {-0.30,-0.15,0,0.15,0.30}
				\draw[thin] (\xGdue+\x, \yG-\thA) -- ++(-0.10, -0.13);
				
				%% ── Anellini ──
				% A_1, A_2: semicerchi (corpo C ruota libero rispetto alla biella)
				\draw[fill=white] 
				(\xAuno-\rn, \yA) arc(180:360:\rn) -- cycle;
				\draw[fill=white] 
				(\xAdue-\rn, \yA) arc(180:360:\rn) -- cycle;
				% B_1, B_2: cerchi pieni (nodi del parallelogramma)
				\draw[fill=white] (\xBuno, \yB) circle (3pt);
				\draw[fill=white] (\xBdue, \yB) circle (3pt);
				% G_1, G_2: cerchi pieni (perni al suolo)
				\draw[fill=white] (\xGuno, \yG) circle (3pt);
				\draw[fill=white] (\xGdue, \yG) circle (3pt);
				
				%				
				%% ── Retta impropria r_∞ (luogo del CR all'infinito) ──
				\def\xRinf{-2.40}
				\def\yRinf{ 1.20}
				\def\angRinf{60}
				\def\Lrinf{1.20}
				\pgfmathsetmacro{\dxR}{\Lrinf*cos(\angRinf)}
				\pgfmathsetmacro{\dyR}{\Lrinf*sin(\angRinf)}
				\draw[gray!60, line width=0.6pt] 
				({\xRinf-\dxR}, {\yRinf-\dyR}) -- ({\xRinf+\dxR}, {\yRinf+\dyR});
				
				%% Etichetta "r_∞" all'estremità superiore della retta
				\node[gray, font=\footnotesize, anchor=south west] 
				at ({\xRinf+\dxR+0.05}, {\yRinf+\dyR+0.05}) {$r_\infty$};
				
				%% Etichetta "C ∈ r_∞" sopra la retta a metà, ruotata
				\node[gray, font=\footnotesize, rotate=\angRinf, anchor=south, 
				inner sep=2pt] 
				at (\xRinf, \yRinf) {$C \in r_\infty$};
				
				
				%% ── Etichetta pannello ──
				\node[anchor=north, font=\small] 
				at (-0.25, \yG-\thA-0.50) {$(a)$};
				
			\end{scope}
			
			% ═══════════════════════════════════════════════════════════
			% PANNELLO (b) — Prestazione statica del bipendolo
			% (realizzazione: pantografo a parallelogramma)
			% ═══════════════════════════════════════════════════════════
			\begin{scope}[xshift=5cm]
				
				%% ── Parametri (identici ad (a)) ──
				%% ── Parametri ──
				\def\xAuno{-0.40}
				\def\xAdue{ 0.40}
				\def\yA{ 0.00}
				\def\Lbiella{0.70}      % era 0.90
				\def\angSup{-110}
				\def\angInf{-70}
				\def\rn{0.10}
				\def\tw{0.22}
				\def\thA{0.32}
				\def\gw{0.40}
				
				% Coordinate calcolate
				\pgfmathsetmacro{\dxSup}{\Lbiella*cos(\angSup)}
				\pgfmathsetmacro{\dySup}{\Lbiella*sin(\angSup)}
				\pgfmathsetmacro{\dxInf}{\Lbiella*cos(\angInf)}
				\pgfmathsetmacro{\dyInf}{\Lbiella*sin(\angInf)}
				\pgfmathsetmacro{\xBuno}{\xAuno+\dxSup}
				\pgfmathsetmacro{\yB}{\yA+\dySup}
				\pgfmathsetmacro{\xBdue}{\xAdue+\dxSup}
				\pgfmathsetmacro{\xGuno}{\xBuno+\dxInf}
				\pgfmathsetmacro{\yG}{\yB+\dyInf}
				\pgfmathsetmacro{\xGdue}{\xBdue+\dxInf}
				
				%% ── Corpo C (identico ad (a)) ──
				\draw[line width=1pt] 
				plot[smooth cycle, tension=0.7] coordinates 
				{(-0.40, 0.00) (-1.30, 0.35) (-1.80, 1.15) 
					(-1.45, 2.05) (-0.55, 2.40) (0.55, 2.40) 
					(1.45, 2.05) (1.80, 1.15) (1.30, 0.35) (0.40, 0.00)};
				\node at (-0.65, 1.65) {$\mathcal{C}$};
				
				
				
				%% ── Bielle superiori, traversa, bielle inferiori (identici ad (a)) ──
				\draw[line width=0.8pt] (\xAuno, \yA) -- (\xBuno, \yB);
				\draw[line width=0.8pt] (\xAdue, \yA) -- (\xBdue, \yB);
				\draw[line width=0.8pt] (\xBuno, \yB) -- (\xBdue, \yB);
				\draw[line width=0.8pt] (\xBuno, \yB) -- (\xGuno, \yG);
				\draw[line width=0.8pt] (\xBdue, \yB) -- (\xGdue, \yG);
				
				%% ── Cerniere al suolo G_1, G_2 (identiche ad (a)) ──
				\draw (\xGuno, \yG) 
				-- ++(-\tw, -\thA) -- ++(2*\tw, 0) -- cycle;
				\draw (\xGuno-\gw, \yG-\thA) -- (\xGuno+\gw, \yG-\thA);
				\foreach \x in {-0.30,-0.15,0,0.15,0.30}
				\draw[thin] (\xGuno+\x, \yG-\thA) -- ++(-0.10, -0.13);
				\draw (\xGdue, \yG) 
				-- ++(-\tw, -\thA) -- ++(2*\tw, 0) -- cycle;
				\draw (\xGdue-\gw, \yG-\thA) -- (\xGdue+\gw, \yG-\thA);
				\foreach \x in {-0.30,-0.15,0,0.15,0.30}
				\draw[thin] (\xGdue+\x, \yG-\thA) -- ++(-0.10, -0.13);
				
				%% ── Anellini (identici ad (a)) ──
				\draw[fill=white] 
				(\xAuno-\rn, \yA) arc(180:360:\rn) -- cycle;
				\draw[fill=white] 
				(\xAdue-\rn, \yA) arc(180:360:\rn) -- cycle;
				\draw[fill=white] (\xBuno, \yB) circle (3pt);
				\draw[fill=white] (\xBdue, \yB) circle (3pt);
				\draw[fill=white] (\xGuno, \yG) circle (3pt);
				\draw[fill=white] (\xGdue, \yG) circle (3pt);
				
				%% ──────────────────────────────────────────────────
				%% COPPIA REATTIVA mu — arco con freccia sopra C
				%% ──────────────────────────────────────────────────
				\def\xMu{0.00}
				\def\yMu{0.50}
				\def\rMu{0.40}
				\draw[->, line width=1.4pt] ({\xMu+\rMu}, \yMu) arc[start angle=0, end angle=180, radius=\rMu];
				\node[font=\footnotesize] at (0.00, 0.5) {$\mu$};
				
				%% ── Etichetta pannello ──
				\node[anchor=north, font=\small] at (-0.25, \yG-\thA-0.50) {$(b)$};
				
			\end{scope}
			
		\end{tikzpicture}
		%	}
	
	\caption[Bipendolo esterno: prestazioni cinematica e statica.]{%
		Bipendolo esterno applicato al corpo $\mathcal{C}$,
		realizzato meccanicamente come pantografo a doppio
		parallelogramma. Prestazione cinematica (a): l'equazione
		di vincolo $\vartheta = 0$ annulla la rotazione del
		corpo, lasciando libere le due componenti di traslazione,
		non rappresentate; il \CRa\ giace sulla retta impropria
		$r_\infty$. Prestazione statica (b): la reazione vincolare
		si riduce alla sola coppia $\mu$, in dualità con la
		rotazione bloccata; la forza risultante è identicamente
		nulla, in accordo con la molteplicità uno.}
	\label{fig:bipendolo_esterno}
\end{figure}
